\scanpage{284}
\begin{center}- 32 -\end{center}\hfill IX

\uline{Proposition} 3.5. (\uline{critère galoisien de réduction semi-stable}). \uline{On suppose} $S$ \uline{un trait}. \uline{Les conditions suivantes sur le schéma abélien} $A_K$ \uline{sont équivalentes}\quad (\uline{les notations étant celles de} (2.5.3), \uline{appliquées à la situation déduite de} $(A_K,S)$ \uline{par passage au hensélisé de} $S$) :

(i) $A_K$ \uline{a réduction semi-stable sur} $S$ (3.4).

(ii) \uline{On a}, \uline{avec les notations} (2.5.3), (2.5.4) :
\[
V\subset V^\perp\quad.
\]
(iii) \uline{On a pour la partie torique} $W$ \uline{du module de Tate} :
\[
W=V^\perp\quad.
\]
(iv) \uline{L'action du groupe d'inertie} $I$ \uline{sur} $U=T_\ell(A(\overline K))$ \uline{est unipotente d'échelon} 2 (\uline{i.e.}\ \uline{il existe un sous-module} $U'$ \uline{de} $U$ \uline{stable par} $I$ \uline{tel que} $I$ \uline{opère trivialement sur} $U'$ \uline{et sur} $U/U'$).

On peut supposer $S$ hensélien (3.4.0). Avec les notations de 2.1.9, (i) signifie que l'on a
\[
\lambda=0\quad,
\]
ce qui en vertu de (2.5.5) équivaut à $V=V\cap V^\perp$ i.e. à (ii). Comme $W=V\cap V^\perp$ par le théorème d'orthogonalité 2.4, on a (ii) $\Longleftrightarrow$ (iii). D'ailleurs (iv) signifie aussi que $I$ opère trivialement sur $U^I\unsure{=|\ne}V$, ce qui par dualité signifie qu'il opère trivialement sur $V^\perp$ i.e. que $V^\perp\subset U^I=V$, ce qui est (ii). Cela achève la démonstration.
\margin{\unsure{q}}

\uline{Remarques} 3.5.1. Les conditions (ii) et (iv) ne dépendent visiblement que de l'action de $I$ sur le $\mathbb Q_\ell$-vectoriel $U\otimes_{\mathbb Z_\ell}\mathbb Q_\ell$, et auraient pu être formulées en termes de cette action. Signalons aussi que ces conditions sont manifestement stables par tout changement de base par un morphisme
\margin{\unsure{?}\quad\unsure{\supset}}

\margin{350}
\margin{350}
\margin{284}
